\documentclass[preview]{standalone}

\usepackage{amsmath}
\usepackage{amssymb}
\usepackage{stellar}
\usepackage{definitions}
\usepackage{bettelini}

\begin{document}

\id{probability-absolute-continuity}
\genpage

\section{Absolute continuity}

\begin{snippetdefinition}{absolute-continuity-lebesgue-definition}{Absolute continuity with respect to Lebesgue measure}
    Let \(\mu\) be a \probabilitymeasure on
    \((\realnumbers,\mathcal B(\realnumbers))\), and let \(\lambda\) denote
    \lebesguemeasure on \(\realnumbers\). The measure \(\mu\) is
    \emph{absolutely continuous with respect to Lebesgue measure}, written
    \[
        \mu\ll\lambda,
    \]
    if, for every \(A\in\mathcal B(\realnumbers)\),
    \[
        \lambda(A)=0
        \quad\Longrightarrow\quad
        \mu(A)=0.
    \]
\end{snippetdefinition}

\begin{snippetdefinition}{probability-density-definition}{Density of an absolutely continuous probability measure}
    Let \(\mu\) be a \probabilitymeasure on
    \((\realnumbers,\mathcal B(\realnumbers))\). A \emph{density} of \(\mu\)
    with respect to \lebesguemeasure is a nonnegative Borel
    \function
    \[
        f_\mu\colon\realnumbers\fromto[0,+\infty]
    \]
    such that
    \[
        \mu(A)=\int_A f_\mu(x)\,dx
    \]
    for every \(A\in\mathcal B(\realnumbers)\).
\end{snippetdefinition}

\begin{snippettheorem}{absolute-continuity-density-characterization}{Density characterization of absolute continuity}
    Let \(\mu\) be a \probabilitymeasure on
    \((\realnumbers,\mathcal B(\realnumbers))\), and let \(\lambda\) be
    \lebesguemeasure. Then
    \[
        \mu\ll\lambda
    \]
    if and only if \(\mu\) admits a \probabilitydensity with respect to
    \(\lambda\).
\end{snippettheorem}

\begin{snippetproof}{absolute-continuity-density-characterization-proof}{absolute-continuity-density-characterization}{Density characterization of absolute continuity}
    If \(\mu(A)=\int_A f_\mu(x)\,dx\) and \(\lambda(A)=0\), then
    \(\int_A f_\mu(x)\,dx=0\), so \(\mu(A)=0\). Conversely, if
    \(\mu\ll\lambda\), the Radon-Nikodym theorem gives a nonnegative Borel
    \function \(f_\mu\) such that
    \[
        \mu(A)=\int_A f_\mu(x)\,dx
    \]
    for every \(A\in\mathcal B(\realnumbers)\).
\end{snippetproof}

\begin{snippetproposition}{probability-density-properties}{Properties of probability densities}
    Let \(\mu\) be an \absolutelycontinuousmeasure on
    \((\realnumbers,\mathcal B(\realnumbers))\), and let \(f_\mu\) be a
    \probabilitydensity of \(\mu\). Then:
    \begin{enumerate}
        \item \(f_\mu\geq0\) almost everywhere;
        \item
            \[
                \int_\realnumbers f_\mu(x)\,dx=1;
            \]
        \item if \(g_\mu\) is another \probabilitydensity of \(\mu\), then
            \(f_\mu=g_\mu\) almost everywhere;
        \item if \(F_\mu\) is the \distributionfunction of \(\mu\), then
            \[
                F_\mu(x)=\int_{-\infty}^x f_\mu(t)\,dt
            \]
            for every \(x\in\realnumbers\), and \(F_\mu\) is continuous.
    \end{enumerate}
\end{snippetproposition}

\begin{snippetproof}{probability-density-properties-proof}{probability-density-properties}{Properties of probability densities}
    Nonnegativity is part of the definition up to modification on a null
    \set. Since \(\mu(\realnumbers)=1\),
    \[
        1=\mu(\realnumbers)=\int_\realnumbers f_\mu(x)\,dx.
    \]
    If \(f_\mu\) and \(g_\mu\) are both densities, then
    \[
        \int_A f_\mu(x)\,dx=\int_A g_\mu(x)\,dx
    \]
    for every Borel \set \(A\). Therefore \(f_\mu=g_\mu\) almost
    everywhere. Finally,
    \[
        F_\mu(x)=\mu((-\infty,x])=\int_{-\infty}^x f_\mu(t)\,dt.
    \]
    Absolute continuity of the Lebesgue integral gives continuity of
    \(F_\mu\).
\end{snippetproof}

\begin{snippettheorem}{probability-measure-from-density}{Probability measure induced by a density}
    Let \(f\colon\realnumbers\fromto[0,+\infty]\) be a Borel \function such
    that
    \[
        \int_\realnumbers f(x)\,dx=1.
    \]
    Then
    \[
        \mu(A)=\int_A f(x)\,dx,
        \qquad A\in\mathcal B(\realnumbers),
    \]
    defines a unique \absolutelycontinuousmeasure on
    \((\realnumbers,\mathcal B(\realnumbers))\), and \(f\) is a
    \probabilitydensity of \(\mu\).
\end{snippettheorem}

\begin{snippetproof}{probability-measure-from-density-proof}{probability-measure-from-density}{Probability measure induced by a density}
    Countable additivity follows from countable additivity of the Lebesgue
    integral on nonnegative functions. Moreover,
    \[
        \mu(\realnumbers)=\int_\realnumbers f(x)\,dx=1.
    \]
    Thus \(\mu\) is a \probabilitymeasure. Its absolute continuity follows
    because the integral of \(f\) over a Lebesgue-null \set is \(0\). The
    formula itself says that \(f\) is a density.
\end{snippetproof}

\section{Distribution functions and densities}

\begin{snippetproposition}{cdf-from-density-and-density-from-cdf}{Distribution function of a density}
    Let \(\mu\) be an \absolutelycontinuousmeasure on
    \((\realnumbers,\mathcal B(\realnumbers))\), let \(f_\mu\) be a
    \probabilitydensity of \(\mu\), and let \(F_\mu\) be its
    \distributionfunction. Then
    \[
        F_\mu(x)=\int_{-\infty}^x f_\mu(t)\,dt
    \]
    for every \(x\in\realnumbers\). Conversely, if a distribution function
    \(F\) has the form
    \[
        F(x)=\int_{-\infty}^x f(t)\,dt
    \]
    for some nonnegative Borel \function \(f\) with
    \(\int_\realnumbers f(t)\,dt=1\), then the corresponding probability
    measure is absolutely continuous with density \(f\).
\end{snippetproposition}

\begin{snippetproof}{cdf-from-density-and-density-from-cdf-proof}{cdf-from-density-and-density-from-cdf}{Distribution function of a density}
    The first statement is the definition of the distribution function
    together with the density formula:
    \[
        F_\mu(x)=\mu((-\infty,x])=\int_{-\infty}^x f_\mu(t)\,dt.
    \]
    The converse follows from the probability measure induced by a density
    and from the fact that a distribution function determines its law.
\end{snippetproof}

\begin{snippetproposition}{piecewise-c1-cdf-density-criterion}{Piecewise differentiable distribution functions}
    Let \(\mu\) be a \probabilitymeasure on
    \((\realnumbers,\mathcal B(\realnumbers))\), and let \(F_\mu\) be its
    \distributionfunction. Suppose that \(F_\mu\) is continuous and that
    there is a finite \set \(S\subseteq\realnumbers\) such that
    \(F_\mu\) is \(\mathcal C^1\) on each connected component of
    \(\realnumbers\difference S\), with
    \[
        F_\mu(y)-F_\mu(x)=\int_x^y F_\mu'(t)\,dt
    \]
    for every \(x<y\). Then \(\mu\) is an \absolutelycontinuousmeasure and
    \[
        f_\mu(x)=F_\mu'(x)
    \]
    almost everywhere is a \probabilitydensity of \(\mu\).
\end{snippetproposition}

\begin{snippetproof}{piecewise-c1-cdf-density-criterion-proof}{piecewise-c1-cdf-density-criterion}{Piecewise differentiable distribution functions}
    Define \(f(x)=F_\mu'(x)\) outside \(S\), and define \(f(x)=0\) on \(S\).
    Since \(F_\mu\) is nondecreasing, \(f\geq0\) almost everywhere. The
    displayed integral identity gives
    \[
        F_\mu(x)=\int_{-\infty}^x f(t)\,dt
    \]
    for every \(x\). Since \(F_\mu\) is a distribution function, this integral
    tends to \(1\) as \(x\to+\infty\). Hence \(f\) is a density, and the law
    determined by \(F_\mu\) is the law induced by \(f\).
\end{snippetproof}

\begin{snippetexample}{continuous-cdf-not-absolutely-continuous-example}{A continuous distribution function need not be absolutely continuous}
    The Cantor distribution has a continuous \distributionfunction, but it is
    singular with respect to \lebesguemeasure. Hence continuity of a
    distribution function does not imply absolute continuity of the
    corresponding probability measure.
\end{snippetexample}

\section{Mixed laws}

\begin{snippetdefinition}{dirac-measure-definition}{Dirac measure}
    Let \(a\in\realnumbers\). The \emph{Dirac measure} at \(a\) is the
    \probabilitymeasure \(\delta_a\) on
    \((\realnumbers,\mathcal B(\realnumbers))\) defined by
    \[
        \delta_a(A)
        =
        \begin{cases}
            1 & a\in A,\\
            0 & a\notin A,
        \end{cases}
    \]
    for every \(A\in\mathcal B(\realnumbers)\).
\end{snippetdefinition}

\begin{snippettheorem}{lebesgue-decomposition-probability-measure}{Lebesgue decomposition of a probability measure}
    Let \(\mu\) be a \probabilitymeasure on
    \((\realnumbers,\mathcal B(\realnumbers))\). Then \(\mu\) can be
    decomposed as
    \[
        \mu=\alpha\mu_d+\beta\mu_{ac}+\gamma\mu_s,
    \]
    where \(\alpha,\beta,\gamma\in[0,1]\), \(\alpha+\beta+\gamma=1\),
    \(\mu_d\) is discrete, \(\mu_{ac}\) is absolutely continuous with respect
    to \lebesguemeasure, and \(\mu_s\) is singular continuous. Terms with
    coefficient \(0\) are omitted.
\end{snippettheorem}

\begin{snippetexample}{mixed-atom-density-lifetime-law}{A mixed lifetime law}
    Let \(p\in[0,1]\) and \(\theta>0\). Let
    \[
        f(x)=\theta e^{-\theta x}\indicator_{[0,+\infty)}(x),
    \]
    and let \(\nu\) be the \probabilitymeasure induced by \(f\). Define
    \[
        \mu=p\delta_0+(1-p)\nu.
    \]
    Then
    \[
        F_\mu(x)
        =
        \begin{cases}
            0 & x<0,\\
            1-(1-p)e^{-\theta x} & x\geq0.
        \end{cases}
    \]
    If \(p>0\), then \(\mu\) is not absolutely continuous with respect to
    \lebesguemeasure, because
    \[
        \mu(\{0\})=p
        \qquad\text{while}\qquad
        \lambda(\{0\})=0.
    \]
    Its non-atomic part on \((0,+\infty)\) has density
    \[
        (1-p)\theta e^{-\theta x}\indicator_{(0,+\infty)}(x).
    \]
\end{snippetexample}

\end{document}
